\documentclass[11pt]{article}
\usepackage{amsmath,amssymb,amsthm}
\usepackage[margin=1in]{geometry}

\newtheorem{definition}{Definition}
\newtheorem{proposition}{Proposition}

\title{Resp\_lc Finite Toy Metric-Response Model}
\author{The AEther-Flow Research Project}
\date{2026-06-18}

\begin{document}
\maketitle

\section*{Control Status}

This artifact is a draft/control Candidate Constructor packet for
`RT-20260614-053`. It constructs a finite toy metric-response model
selected by `RT-20260614-052`. It is not canonical ontology, not a
canonical `Resp\_lc` witness, not `M\_src`, not `g\_eff`, not matter
coupling, not Einstein equations, not benchmark evidence, not a Gate
Chair verdict, and not completed-derivation evidence.

\section{Finite Source Set}

Let the toy cell set be
\[
  C_{\mathrm{toy}}=\{U_0,U_1\}.
\]
Let the response-token alphabet be
\[
  T_{\mathrm{toy}}=\{a,b\}.
\]
A tagged toy source object is a tuple
\[
  X=(C_{\mathrm{toy}},E,\varepsilon,\lambda,\tau),
\]
where \(E=\{\{U_0,U_1\}\}\), \(\varepsilon\in\{+1,-1\}\) is an
explicit source-side orientation tag, \(\lambda\in\{1,2\}\) is an
explicit positive source-side normalization tag, and
\(\tau:C_{\mathrm{toy}}\rightarrow T_{\mathrm{toy}}\) is an explicit
source-side token tag.

The finite family is
\[
\begin{aligned}
F_{\mathrm{toy}}=\{&
X_+       =(C_{\mathrm{toy}},E,+1,1,\tau_{ab}),\\
&X_{\mathrm{flip}}=(C_{\mathrm{toy}},E,-1,1,\tau_{ab}),\\
&X_{\mathrm{scale}}=(C_{\mathrm{toy}},E,+1,2,\tau_{ab}),\\
&X_{\mathrm{token}}=(C_{\mathrm{toy}},E,+1,1,\tau_{ba}),\\
&X_{\emptyset}=(C_{\mathrm{toy}},E)\}.
\end{aligned}
\]
Here \(\tau_{ab}(U_0)=a\), \(\tau_{ab}(U_1)=b\),
\(\tau_{ba}(U_0)=b\), and \(\tau_{ba}(U_1)=a\). The object
\(X_{\emptyset}\) deliberately has no orientation, normalization, or
token tags. It models the current obstruction: the retained tuple by
itself does not determine the response data.

\section{Partial Response Relation}

\begin{definition}[Toy response relation]
Define a partial relation
\[
  R_{\mathrm{toy}}:F_{\mathrm{toy}}\rightharpoonup
  \{+1,-1\}\times\{1,2\}\times T_{\mathrm{toy}}^{C_{\mathrm{toy}}}
\]
by
\[
  R_{\mathrm{toy}}(C_{\mathrm{toy}},E,\varepsilon,\lambda,\tau)
  =(\varepsilon,\lambda,\tau)
\]
whenever all three tags are explicit source data. The relation is
undefined on \(X_{\emptyset}\).
\end{definition}

This relation is intentionally conservative. It does not choose a sign
by convention, it does not choose a scale by calibration, and it does
not assign response tokens by notation. It records only tags already
present in the finite toy source object.

\section{Toy Metric-Response Analogue}

\begin{definition}[Toy distance and signed adjacency]
For each tagged object \(X\) with
\(R_{\mathrm{toy}}(X)=(\varepsilon,\lambda,\tau)\), define
\[
  D_X(U_i,U_j)=
  \begin{cases}
    0,& i=j,\\
    \lambda,& i\neq j,
  \end{cases}
\]
and define the signed adjacency form
\[
  A_X(U_i,U_j)=
  \begin{cases}
    0,& i=j,\\
    \varepsilon\lambda,& i\neq j.
  \end{cases}
\]
The token component is the explicit map \(\tau\).
\end{definition}

\(D_X\) and \(A_X\) are finite response analogues only. They are not a
Lorentzian metric, not `M\_src`, not `g\_eff`, and not a matter
coupling law.

\section{Invariance Checks}

\begin{proposition}[Finite toy covariance]
The finite construction has the following controlled invariance
properties.
\end{proposition}

\begin{proof}
First, let \(\pi\) be the nontrivial permutation of
\(C_{\mathrm{toy}}\). Relabeling cells sends \(D_X\) and \(A_X\) to
their conjugates under \(\pi\). The numerical orientation and
normalization tags are unchanged, and the token map is transported by
precomposition with \(\pi^{-1}\). Thus the construction is covariant
under finite cell relabeling.

Second, let \(\sigma\) be a permutation of the token alphabet
\(T_{\mathrm{toy}}\). Replacing \(\tau\) by \(\sigma\circ\tau\)
changes only the token names. The orientation tag \(\varepsilon\), the
normalization tag \(\lambda\), \(D_X\), and \(A_X\) are unchanged.

Third, external sign flips or external rescalings that are not source
tags do not change \(R_{\mathrm{toy}}\). The relation reads only
\(\varepsilon\), \(\lambda\), and \(\tau\) as explicit toy source data.

Fourth, finite tag perturbations change only the matching response
component. \(X_{\mathrm{flip}}\) changes the signed adjacency sign,
\(X_{\mathrm{scale}}\) changes the off-diagonal distance weight, and
\(X_{\mathrm{token}}\) changes only token assignment. Removing the tags
gives \(X_{\emptyset}\), where \(R_{\mathrm{toy}}\) is undefined.
\end{proof}

\section{Bridge Attempt Status}

The finite candidate map is
\[
  X\longmapsto (R_{\mathrm{toy}}(X),D_X,A_X,\tau)
\]
for explicitly tagged \(X\in F_{\mathrm{toy}}\). The map succeeds only
as a finite toy construction. It does not solve the missing primitive
for the current untagged retained tuple. That missing primitive remains
an actual source-side response law or admissible source extension that
determines orientation sign, positive normalization, and response-token
semantics without target-GR import.

\section{Distance-to-GR Status}

\begin{center}
\begin{tabular}{p{0.33\linewidth}p{0.57\linewidth}}
\hline
Burden & Status \\
\hline
Source ontology primitives & Unchanged; no canonical ontology edit is made. \\
Source equivalence EqSrc & Unchanged; general EqSrc remains undischarged. \\
RetainH & Unchanged; no canonical RetainH adoption is supplied. \\
GenH & Unchanged; no canonical GenH adoption is supplied. \\
ObsLoc\_lc & Unchanged; local exact-branch witness status is not promoted. \\
Resp\_lc & Still blocked for the current untagged tuple; toy tags do not adopt Resp\_lc. \\
M\_src & Not started; no source manifold construction is supplied. \\
g\_eff & Not started; no effective metric is supplied. \\
Matter coupling & Not started. \\
Einstein equations & Not started. \\
Finite-variation robustness & Not discharged beyond finite toy tag checks. \\
Benchmark promotion & Blocked. \\
Gate Chair status & Blocked and human-gated. \\
Current route freeze or hard-fail status & Not frozen; finite toy construction now requires audit. \\
\hline
\end{tabular}
\end{center}

\section{Conclusion}

The finite toy metric-response model exists as draft/control
scaffolding. It demonstrates that a response relation can be made
source-determined in a toy setting when orientation, normalization, and
token semantics are explicit source tags. It also preserves the core
obstruction: without those tags, the relation is undefined. The logical
next step is a bounded Smuggling Auditor audit for target-GR import,
arbitrary convention leakage, and ontology-overread.

\begin{thebibliography}{9}

\bibitem{burden-map}
The AEther-Flow Research Project. (2026, June 17). \emph{GR derivation
burden map} [Internal control note].

\bibitem{handoff0094}
The AEther-Flow Research Project. (2026, June 17). \emph{Handoff 0094}
[Internal research-control handoff].

\bibitem{selector-finite-toy}
The AEther-Flow Research Project. (2026, June 17). \emph{Resp\_lc
theoretical continuation selector finite toy decision} [Internal
research-control artifact].

\end{thebibliography}

\end{document}
